\documentclass[twocolumn,11pt,a4paper]{article}

% ============================================================
%  TRANSACTIONAL MAGNITUDE CALCULUS:
%  MONETARY DERIVATIVES, S-ENTROPY NORMALISATION,
%  AND THE MONETARY TANGENT SPACE
% ============================================================

\usepackage[utf8]{inputenc}
\usepackage[T1]{fontenc}
\usepackage{lmodern}
\usepackage[a4paper,margin=2.5cm]{geometry}
\usepackage{amsmath,amssymb,amsthm,mathtools}
\usepackage{bm}
\usepackage{bbm}
\usepackage{booktabs}
\usepackage{array}
\usepackage{enumitem}
\usepackage{hyperref}
\usepackage{xcolor}
\usepackage[round]{natbib}
\usepackage{microtype}

\hypersetup{
  colorlinks=true,
  linkcolor=blue!40!black,
  citecolor=blue!40!black,
  urlcolor=blue!40!black
}

% Theorem environments
\newtheoremstyle{thmstyle}
  {\topsep}{\topsep}
  {\itshape}{}
  {\bfseries}{.}
  {.5em}{}
\theoremstyle{thmstyle}
\newtheorem{theorem}{Theorem}[section]
\newtheorem{lemma}[theorem]{Lemma}
\newtheorem{proposition}[theorem]{Proposition}
\newtheorem{corollary}[theorem]{Corollary}

\theoremstyle{definition}
\newtheorem{definition}[theorem]{Definition}
\newtheorem{example}[theorem]{Example}
\newtheorem{construction}[theorem]{Construction}

\theoremstyle{remark}
\newtheorem{remark}[theorem]{Remark}
\newtheorem{notation}[theorem]{Notation}

% Core operators
\DeclareMathOperator{\Var}{Var}
\DeclareMathOperator{\Cov}{Cov}
\DeclareMathOperator{\esssup}{ess\,sup}
\DeclareMathOperator{\sgn}{sgn}

% Monetary calculus macros
\newcommand{\Gain}{G}                        % gain-loss process G(t,\tau)
\newcommand{\Clock}{\Theta}                  % transaction clock \Theta(t)
\newcommand{\InvClock}{\theta}               % right-continuous inverse \theta(s)
\newcommand{\Mdiff}[1]{\frac{d#1}{d\Clock}} % monetary derivative d/d\Theta
\newcommand{\MdiffK}[2]{\frac{d#1}{d\Clock_{#2}}} % layer-k monetary derivative
\newcommand{\Sval}{S}                        % S-functional value
\newcommand{\Sscale}{\mathbb{S}}             % S-scale [0,100]
\newcommand{\Sfloor}{S_{\flat}}              % floor function
\newcommand{\Mtan}{\mathcal{T}^{\Clock}}     % monetary tangent space
\newcommand{\wstar}{w^{*}}                   % fixed-point portfolio
\newcommand{\Lop}{\mathbf{L}}               % graph Laplacian
\newcommand{\Lpinv}{\mathbf{L}^{\dagger}}   % pseudoinverse
\newcommand{\lamtwo}{\lambda_{2}}           % Fiedler value
\newcommand{\Deltam}{\Delta_{m}}            % simplex
\newcommand{\Real}{\mathbb{R}}
\newcommand{\Prob}{\mathbb{P}}
\newcommand{\Expect}{\mathbb{E}}
\newcommand{\Filtr}{\mathcal{F}}
\newcommand{\Borel}{\mathcal{B}}
\newcommand{\Receiver}{\mathcal{R}}

\title{\Large\textbf{Transactional Magnitude Calculus:\\
Monetary Derivatives, S-Entropy Normalisation,\\
and the Monetary Tangent Space}}

\author{\textit{Anonymous submission for review}}

\date{}

\begin{document}
\maketitle

% ============================================================
\begin{abstract}
Classical analysis measures rates of change relative to
calendar time.  In systems driven by economic activity,
however, the natural clock is not the uniform tick of the
calendar but the accumulated magnitude of transactional
events.  We develop a rigorous calculus in which the
fundamental rate operator is the \emph{monetary derivative}
$d/d\Clock$, where $\Clock(t) = \int_{0}^{t}|\Gain(s)|\,ds$
is the \emph{transaction clock}---the running integral of
absolute gain-loss intensity.  Replacing calendar time with
transaction time subordinates every process to the
underlying economic activity and eliminates the
dimensionality barriers that obstruct cross-quantity
comparison: when the S-entropy functional
$\Sval:\Receiver\times\mathcal{X}\times\mathcal{X}^{*}
\to[0,100]$ is used as the scalar field, every monetary
derivative $d\Sval/d\Clock$ is measured in
\emph{S-units per S-unit}---a pure number that adds
across heterogeneous quantities without conversion factors.
We formalise the \emph{monetary tangent space} $\Mtan$ at
each transaction-time point, prove that it carries a natural
norm inherited from the S-functional, and show that it is
closed under the monetary derivative operator.  For systems
governed by a hierarchical gear network of nested stopping
times $\{\sigma_{k}\}_{k\geq 1}$, we derive a
\emph{telescoping gear-ratio formula}
$d\Sval/d\Clock_{K} = (d\Sval/d\Clock_{1})
\prod_{k=1}^{K-1}(\Clock_{k}/\Clock_{k+1})$
and prove that the No-Privileged-Level principle---the same
algebraic laws hold at every gear depth---is an exact
consequence of the S-entropy axioms.  We further prove an
\emph{ergodic consistency theorem}: under a stationary
ergodic return process the time-averaged monetary
derivative converges almost surely to the stationary
mean, and the convergence rate is $O(K^{-1/2})$ by
Birkhoff's theorem applied in transaction time.  The
framework provides a principled foundation for comparing
rates of change of price, volume, volatility, and any
other economic observable on a single dimensionless scale,
and reduces to classical calculus when transaction activity
is constant.
\end{abstract}

\noindent\textbf{Keywords:} transaction clock, monetary
derivative, S-entropy, monetary tangent space, gear network,
No-Privileged-Level, Clark subordination, ergodic theory,
dimensional analysis, gain-loss intensity.

\vspace{6pt}
\tableofcontents
\vspace{6pt}


% ============================================================
\section{Introduction}
\label{sec:intro}
% ============================================================

The calculus of variations, Newtonian mechanics, and
classical stochastic analysis all share a common
primitive: the calendar-time derivative $d/dt$.  For
physical systems whose natural clocking is the uniform
passage of time this is irreproachable.  But financial
markets, biological signalling cascades, supply networks,
and social contagion all possess an internal clock whose
rate is itself driven by the intensity of the phenomenon
being studied.  A stock that trades ten thousand shares in
a second has advanced further along its economic trajectory
than a stock that trades ten shares over a day, even if
both share the same calendar elapsed time.

This observation has antecedents.  Clark~\citep{Clark1973}
showed empirically that daily price changes are better
described by a subordinated process: Brownian motion in
transaction time, where transaction time is the
accumulated volume or trade count.  Mandelbrot and Taylor
\citep{MandelbrotTaylor1967} and
Monroe~\citep{Monroe1978} formalised the connection between
continuous local martingales and Brownian motion in a
subordinated clock.  Ane and Geman~\citep{AneGeman2000}
confirmed empirically that conditioning on trade count
renders intraday returns approximately Gaussian.
In the mathematical finance literature, stochastic time
changes appear in the context of variance processes and
activity rates~\citep{BarndorffNielsenShephard2001}.

Yet despite these advances, no systematic calculus has been
developed that:
\begin{enumerate}[nosep]
\item replaces the calendar-time derivative $d/dt$ with a
      \emph{monetary derivative} $d/d\Clock$ uniformly
      across all quantities of interest;
\item resolves the \emph{dimensional incommensurability}
      problem that arises when one attempts to add rates of
      change of heterogeneous observables (price velocity
      cannot simply be added to volume velocity);
\item provides a \emph{scale-free, addable} scalar for
      rates of change that respects the algebraic structure
      of bounded-cognition information theory.
\end{enumerate}

The present paper addresses all three gaps.  The key
insight is that when rates of change are measured in the
S-entropy scale $[0,100]$ of
\citet{Sachikonye2024S}---a bounded, receiver-relative
distance-to-truth functional---and normalised by the
transaction clock, the result is a pure number: S-units of
distance-to-truth per S-unit of distance-to-truth.
Two monetary derivatives of entirely different quantities
are then directly addable, subtractable, and comparable.

The paper is organised as follows.  Section~\ref{sec:prelim}
sets up notation and reviews the required background on
graph Laplacians, Kirchhoff portfolios, and S-entropy.
Section~\ref{sec:clock} defines the transaction clock and
proves its basic properties.  Section~\ref{sec:mderiv}
introduces the monetary derivative operator and develops
its analysis.  Section~\ref{sec:sentropy} shows how
S-entropy eliminates dimensional barriers.
Section~\ref{sec:tangent} constructs the monetary tangent
space and establishes its normed-space structure.
Section~\ref{sec:gear} derives the telescoping gear-ratio
formula for hierarchical networks.  Section~\ref{sec:npl}
proves the No-Privileged-Level theorem in transaction time.
Section~\ref{sec:ergodic} establishes ergodic consistency
of monetary derivatives.  Section~\ref{sec:examples}
works three illustrative examples.  Section~\ref{sec:conclude}
discusses implications and open questions.


% ============================================================
\section{Preliminaries}
\label{sec:prelim}
% ============================================================

\subsection{Probability and filtration}

Let $(\Omega,\Filtr,\Prob)$ be a complete probability space
and $(\Filtr_{t})_{t\geq 0}$ a right-continuous filtration
satisfying the usual conditions.  All processes are assumed
to be $(\Filtr_{t})$-adapted.

\subsection{Graph Laplacian and Kirchhoff portfolio}

Let $G=(V,E,A)$ be a weighted undirected graph with
$m=|V|$ assets, adjacency matrix $A\in\Real^{m\times m}$
($A_{ij}=A_{ji}\geq 0$, $A_{ii}=0$), degree matrix
$D=\mathrm{diag}(A\mathbf{1})$, and Laplacian
$\Lop=D-A$.  Since $G$ is connected, $\Lop$ has
eigenvalues $0=\lambda_{1}<\lamtwo\leq\cdots\leq\lambda_{m}$;
$\lamtwo$ is the Fiedler (algebraic connectivity) value.

For a forecast return vector $\hat\mu\in\Real^{m}$ with
mean-centred version $\mu_{c}=\hat\mu-\bar{\hat\mu}\mathbf{1}$,
the \emph{Kirchhoff fixed-point portfolio} is
\begin{equation}\label{eq:kfp}
  \wstar = \Lpinv\mu_{c} + \tfrac{1}{m}\mathbf{1},
\end{equation}
where $\Lpinv$ is the Moore--Penrose pseudoinverse of
$\Lop$~\citep{Penrose1955}.  The weight $\wstar$ satisfies
the discrete Kirchhoff law $\Lop\wstar=\mu_{c}$ and lies
on the probability simplex $\Deltam$~\citep{Kirchhoff1847}.

\subsection{Gain-loss process}

Let $R(t,\tau)\in\Real^{m}$ be the vector of asset returns
over the interval $[t,t+\tau]$, and let $\hat\mu(\tau)
=\Expect[R(t,t+\tau)]$ be the corresponding forecast.
The \emph{Kirchhoff gain-loss} at calendar time $t$ with
horizon $\tau$ is
\begin{equation}\label{eq:gain}
  \Gain(t,\tau) = (\wstar(\tau))^{\top}
                  [R(t,t+\tau)-\hat\mu(\tau)].
\end{equation}
When the forecast is conditionally unbiased---i.e.\
$\Expect[R(t,t+\tau)\mid\Filtr_{t}]=\hat\mu(\tau)$---the
process $(\Gain(t,\tau))_{t}$ is a martingale difference
sequence with respect to $(\Filtr_{t})$.

\subsection{S-entropy functional}

We recall the S-entropy framework of \citet{Sachikonye2024S}.
Let $\mathcal{X}$ be a candidate space, $\mathcal{X}^{*}
\subseteq\mathcal{X}$ the truth space, and $\Receiver$ the
class of all receivers.  An \emph{S-functional} is a map
\begin{equation}\label{eq:sfunc}
  \Sval : \Receiver \times \mathcal{X} \times \mathcal{X}^{*}
          \to [0,100],
  \quad (\mathcal{R},x,x^{*}) \mapsto \Sval_{\mathcal{R}}(x,x^{*}),
\end{equation}
satisfying non-negativity, boundedness, the positive floor
$\Sfloor(\mathcal{R})>0$ for every bounded receiver
$\mathcal{R}$, and receiver-decoded coherence.  The image
of $\Sval_{\mathcal{R}}$ lies in $[\Sfloor(\mathcal{R}),100]$.

A key consequence---the \emph{Unconstrained Subtask
Theorem} of \citet{Sachikonye2024S}---states that any
S-value $\Sval^{*}$ admits an arbitrary type-mixed
decomposition into sub-expressions, signed or fractional,
provided their composition under the receiver's evaluation
map yields $\Sval^{*}$.  This freedom is what permits
heterogeneous monetary derivatives to be decomposed and
recombined without constraint.


% ============================================================
\section{The Transaction Clock}
\label{sec:clock}
% ============================================================

\begin{definition}[Transaction clock]\label{def:clock}
Let $\Gain(\cdot,\tau)$ be the gain-loss process of
\eqref{eq:gain}.  The \emph{transaction clock} is the
right-continuous, non-decreasing process
\begin{equation}\label{eq:clock}
  \Clock(t) \;:=\; \int_{0}^{t} |\Gain(s,\tau)|\,ds,
  \quad t\geq 0.
\end{equation}
\end{definition}

The transaction clock accumulates the absolute magnitude
of economic activity.  When activity is high, $\Clock$
advances rapidly; when markets are inactive, $\Clock$
nearly stalls.

\begin{lemma}[Basic properties of $\Clock$]\label{lem:clock-props}
Under the standing assumptions:
\begin{enumerate}[nosep,label=(\roman*)]
\item $\Clock(0)=0$ and $\Clock$ is $(\Filtr_{t})$-adapted;
\item $\Clock$ is non-decreasing and absolutely continuous
      with $\dot\Clock(t)=|\Gain(t,\tau)|$ a.e.;
\item $\Clock$ is a finite-variation process with quadratic
      variation $[\Clock]_{t}=0$;
\item if $\Expect[|\Gain(t,\tau)|]<\infty$ for all $t$,
      then $\Expect[\Clock(t)]<\infty$.
\end{enumerate}
\end{lemma}

\begin{proof}
(i) follows from $\Gain$ being adapted and the integral
being a progressively measurable operation.  (ii) is
immediate from the fundamental theorem of calculus for
absolutely continuous functions.  (iii) follows because
$\Clock$ has bounded variation, so its quadratic variation
vanishes.  (iv) follows by Fubini's theorem:
$\Expect[\Clock(t)]=\int_{0}^{t}\Expect[|\Gain(s,\tau)|]\,ds<\infty$.
\end{proof}

\begin{definition}[Transaction-time inverse]\label{def:inv-clock}
The \emph{transaction-time inverse} is the right-continuous
process
\begin{equation}\label{eq:inv-clock}
  \InvClock(s) \;:=\; \inf\{t\geq 0 : \Clock(t) > s\},
  \quad s\geq 0.
\end{equation}
\end{definition}

The pair $(\Clock, \InvClock)$ implements a random time
change in the sense of Clark~\citep{Clark1973}: a process
$X$ observed in calendar time is transformed to
$X_{\InvClock(s)}$ observed in transaction time $s$.

\begin{proposition}[Subordination]\label{prop:subordination}
Let $X(t)$ be a continuous semimartingale.  Then
$Y(s):=X(\InvClock(s))$ is a process in transaction time
whose infinitesimal increments satisfy
\begin{equation}\label{eq:subord}
  dY(s) = dX(\InvClock(s)) \cdot
          \frac{1}{|\Gain(\InvClock(s),\tau)|}\,ds.
\end{equation}
In particular, $dY(s)/ds$ is the rate of change of $X$
per unit of transaction activity.
\end{proposition}

\begin{proof}
By the chain rule for absolutely continuous time changes
\citep{KaratzasShreve1991}, since $\Clock$ is absolutely
continuous with density $|\Gain|$,
$d\InvClock(s)/ds = 1/|\Gain(\InvClock(s))|$.  Therefore
$dY(s) = (dX/dt)|_{t=\InvClock(s)}\cdot
(d\InvClock/ds)\,ds
= (dX/dt)/|\Gain|\,ds$ as claimed.
\end{proof}

\begin{remark}\label{rem:clark-connection}
Proposition~\ref{prop:subordination} generalises Clark's
result from trade-count subordination to gain-loss
subordination.  The economic interpretation is that one
unit of transaction time represents one unit of cumulative
gain-loss magnitude, making the clock a direct measure of
economic intensity rather than a proxy (volume, trade count).
\end{remark}

\begin{theorem}[Variance rescaling under subordination]\label{thm:var-rescale}
Let $X(t)$ be a zero-mean martingale with
$\Var[X(t)]=\sigma^{2}t$ and let $Y(s)=X(\InvClock(s))$.
Then
\begin{equation}\label{eq:var-rescale}
  \Var[Y(s)] = \sigma^{2}\,\Expect[\InvClock(s)].
\end{equation}
Under the ergodic assumption
$\Expect[|\Gain(t)|]\to\bar{g}>0$ as $t\to\infty$,
$\Expect[\InvClock(s)]\sim s/\bar{g}$ and
$\Var[Y(s)]\sim \sigma^{2}s/\bar{g}$.
\end{theorem}

\begin{proof}
By the optional stopping theorem and the tower property,
$\Expect[X(\InvClock(s))^{2}]
= \sigma^{2}\Expect[\InvClock(s)]$, which gives
\eqref{eq:var-rescale}.  The ergodic estimate follows
from the law of large numbers:
$\Clock(t)/t\to\bar{g}$ a.s.\ implies
$\InvClock(s)/s\to 1/\bar{g}$ a.s.\ by the inverse
function theorem for monotone processes, so
$\Expect[\InvClock(s)]\to s/\bar{g}$ by dominated
convergence under integrability conditions.
\end{proof}


% ============================================================
\section{The Monetary Derivative Operator}
\label{sec:mderiv}
% ============================================================

\begin{definition}[Monetary derivative]\label{def:mderiv}
Let $f:\Real_{\geq 0}\to\Real$ be a differentiable function
of calendar time.  Its \emph{monetary derivative} at time
$t$ (with respect to the transaction clock $\Clock$) is
\begin{equation}\label{eq:mderiv}
  \Mdiff{f}(t)
  \;:=\; \frac{df/dt}{|\Gain(t,\tau)|}
  \;=\; \frac{\dot{f}(t)}{|\Gain(t,\tau)|}
\end{equation}
whenever $\Gain(t,\tau)\neq 0$.
\end{definition}

\begin{remark}\label{rem:mderiv-units}
If $f$ is measured in units $[f]$ and calendar time is
measured in seconds, then $df/dt$ has units $[f]\cdot
\mathrm{s}^{-1}$.  The gain-loss $|\Gain|$ is a
dimensionless portfolio return, hence
$\Mdiff{f}$ has units $[f]$ per unit return---a quantity
that still depends on the measurement scale of $f$.
Section~\ref{sec:sentropy} resolves this residual
dimensionality by passing to S-entropy coordinates.
\end{remark}

\begin{lemma}[Linearity]\label{lem:linearity}
For differentiable $f,g$ and constants $\alpha,\beta\in\Real$:
\begin{equation}
  \Mdiff{(\alpha f+\beta g)}
  = \alpha\,\Mdiff{f} + \beta\,\Mdiff{g}.
\end{equation}
\end{lemma}

\begin{proof}
$d(\alpha f+\beta g)/dt = \alpha\dot{f}+\beta\dot{g}$;
divide by $|\Gain|\neq 0$.
\end{proof}

\begin{lemma}[Product rule]\label{lem:product}
For differentiable $f,g$:
\begin{equation}
  \Mdiff{(fg)} = g\,\Mdiff{f} + f\,\Mdiff{g}.
\end{equation}
\end{lemma}

\begin{proof}
$(fg)^{\cdot} = \dot{f}g+f\dot{g}$; divide by $|\Gain|$
and use Definition~\ref{def:mderiv}.
\end{proof}

\begin{lemma}[Chain rule]\label{lem:chain}
Let $f = h \circ g$ where $g:\Real_{\geq 0}\to\Real$
and $h:\Real\to\Real$ are differentiable.  Then
\begin{equation}
  \Mdiff{f}(t) = h'(g(t))\,\Mdiff{g}(t).
\end{equation}
\end{lemma}

\begin{proof}
$\dot{f} = h'(g)\dot{g}$; divide by $|\Gain|$.
\end{proof}

\begin{proposition}[Monetary derivative in transaction time]\label{prop:mderiv-transtime}
Let $Y(s) = f(\InvClock(s))$ be $f$ observed in
transaction time.  Then
\begin{equation}\label{eq:mderiv-transtime}
  \frac{dY}{ds}(s) = \Mdiff{f}(\InvClock(s)).
\end{equation}
That is, the monetary derivative in calendar time
coincides with the ordinary derivative in transaction time.
\end{proposition}

\begin{proof}
$dY/ds = (df/dt)|_{t=\InvClock(s)}\cdot d\InvClock/ds
= \dot{f}(\InvClock(s))/|\Gain(\InvClock(s))|
= \Mdiff{f}(\InvClock(s))$.
\end{proof}

\begin{remark}\label{rem:unification}
Proposition~\ref{prop:mderiv-transtime} shows that working
with monetary derivatives in calendar time is \emph{identical}
to working with ordinary derivatives in transaction time.
The two representations are related by the random time
change $(\Clock,\InvClock)$ and are freely interchangeable.
\end{remark}

\begin{theorem}[Bounded activity bound]\label{thm:bounded-activity}
Suppose $|\dot{f}(t)|\leq C$ for all $t$ and
$|\Gain(t)|\geq g_{\min}>0$ a.s.\ for all $t$.  Then
\begin{equation}
  \left|\Mdiff{f}(t)\right| \leq \frac{C}{g_{\min}}\quad a.s.
\end{equation}
\end{theorem}

\begin{proof}
$|\Mdiff{f}| = |\dot{f}|/|\Gain| \leq C/g_{\min}$.
\end{proof}

\begin{theorem}[Monetary fundamental theorem]\label{thm:mfundamental}
For any differentiable $f$ with integrable $\Mdiff{f}$:
\begin{equation}\label{eq:mfundamental}
  f(t_{2}) - f(t_{1})
  = \int_{\Clock(t_{1})}^{\Clock(t_{2})}
    \Mdiff{f}(\InvClock(s))\,ds.
\end{equation}
\end{theorem}

\begin{proof}
By the change-of-variables formula for monotone increasing
$\Clock$:
$\int_{t_{1}}^{t_{2}} \dot{f}(t)\,dt
= \int_{\Clock(t_{1})}^{\Clock(t_{2})}
  \dot{f}(\InvClock(s))\,\frac{d\InvClock}{ds}\,ds
= \int_{\Clock(t_{1})}^{\Clock(t_{2})}
  \frac{\dot{f}(\InvClock(s))}{|\Gain(\InvClock(s))|}\,ds$.
The result follows from
$f(t_{2})-f(t_{1})=\int_{t_{1}}^{t_{2}}\dot{f}(t)\,dt$.
\end{proof}


% ============================================================
\section{S-Entropy as Dimensionless Currency}
\label{sec:sentropy}
% ============================================================

The monetary derivative $\Mdiff{f}$ in Definition~\ref{def:mderiv}
has residual units $[f]$ inherited from $f$.  Comparing
monetary derivatives of price (in currency units) with
those of volume (in share counts) or volatility (dimensionless)
requires a common unit.  The S-entropy functional
\eqref{eq:sfunc} provides exactly this.

\begin{definition}[S-transformed quantity]\label{def:s-transform}
Let $f:\Real_{\geq 0}\to\Real$ be a quantity of interest
with truth value $f^{*}(t)$ at time $t$, and let
$\mathcal{R}$ be a fixed receiver.  The
\emph{S-transformed quantity} is
\begin{equation}
  \tilde{f}(t)
  := \Sval_{\mathcal{R}}(f(t), f^{*}(t)) \in [0,100].
\end{equation}
\end{definition}

\begin{proposition}[Dimensionlessness of S-transform]
\label{prop:dimensionless}
For any quantities $f$ and $g$ measured in different units,
the S-transforms $\tilde{f}$ and $\tilde{g}$ are both
measured in S-units on the scale $[0,100]$, and their
monetary derivatives $\Mdiff{\tilde{f}}$ and
$\Mdiff{\tilde{g}}$ are both in S-units per unit
transaction magnitude---the \emph{same unit}---making
them directly addable.
\end{proposition}

\begin{proof}
$\tilde{f}$ and $\tilde{g}$ are both images of the
S-functional, which maps into $[0,100]\subset\Real$
with no physical dimension.  The transaction clock $\Clock$
accumulates the dimensionless portfolio return $|\Gain|$,
so $\Clock$ is also dimensionless.  Therefore
$d\tilde{f}/d\Clock$ and $d\tilde{g}/d\Clock$ are both
dimensionless ratios in $\Real$.
\end{proof}

\begin{corollary}[Additivity across heterogeneous quantities]
\label{cor:additivity}
For any finite collection of quantities $f_{1},\ldots,f_{n}$
measured in arbitrary units and any weights
$\alpha_{i}\in\Real$ with $\sum_{i}|\alpha_{i}|<\infty$:
\begin{equation}\label{eq:additivity}
  \sum_{i=1}^{n} \alpha_{i}\,\Mdiff{\tilde{f}_{i}}(t)
  \in \Real
\end{equation}
is well-defined and carries the single interpretation of
a \emph{net monetary derivative in S-units per transaction unit}.
\end{corollary}

\begin{proof}
Each $\Mdiff{\tilde{f}_{i}}$ is a real number by
Proposition~\ref{prop:dimensionless}, and a finite linear
combination of real numbers is a real number.
\end{proof}

\begin{remark}\label{rem:adding-apples-oranges}
Corollary~\ref{cor:additivity} formalises the informal
observation that S-entropy allows one to \emph{add apples
and oranges}: by mapping every quantity into the common
currency $[0,100]$ before differentiating, the monetary
derivative strips away all legacy units and returns a
universal scalar.  The Unconstrained Subtask Theorem of
\citet{Sachikonye2024S} further guarantees that any such
sum can be decomposed into sub-tasks in an arbitrary
type-mixed fashion without violating the total S-value.
\end{remark}

\begin{theorem}[Floor persistence under monetary differentiation]
\label{thm:floor-persistence}
Let $\mathcal{R}$ be a bounded receiver with floor
$\Sfloor(\mathcal{R})>0$.  Then for every quantity $f$
with $\tilde{f}(t)\geq\Sfloor(\mathcal{R})$ for all $t$
and every interval $[t_{1},t_{2}]$:
\begin{equation}
  \tilde{f}(t_{2}) - \tilde{f}(t_{1})
  = \int_{\Clock(t_{1})}^{\Clock(t_{2})}
    \Mdiff{\tilde{f}}(\InvClock(s))\,ds,
\end{equation}
and the integral is bounded:
$|\tilde{f}(t_{2})-\tilde{f}(t_{1})|\leq 100-\Sfloor(\mathcal{R})$.
\end{theorem}

\begin{proof}
The identity follows from Theorem~\ref{thm:mfundamental}
applied to $\tilde{f}$.  The bound follows from
$\tilde{f}\in[\Sfloor(\mathcal{R}),100]$, so
$|\tilde{f}(t_{2})-\tilde{f}(t_{1})|
\leq 100-\Sfloor(\mathcal{R})$.
\end{proof}


% ============================================================
\section{The Monetary Tangent Space}
\label{sec:tangent}
% ============================================================

\begin{definition}[Monetary tangent space]\label{def:tangent}
Fix a receiver $\mathcal{R}$ and a transaction-time point
$s\geq 0$.  The \emph{monetary tangent space at $s$} is
the vector space
\begin{equation}
  \Mtan_{s} := \left\{
    v = \Mdiff{\tilde{f}}(\InvClock(s))
    \;\Big|\;
    f:\Real_{\geq 0}\to\Real\ \text{differentiable},\
    \tilde{f}\ \text{as in Definition~\ref{def:s-transform}}
  \right\} \subseteq \Real.
\end{equation}
\end{definition}

\begin{remark}\label{rem:tangent-is-real}
Since all S-transformed quantities take values in $[0,100]$,
the monetary tangent space $\Mtan_{s}$ is a subspace of
$\Real$.  In the finite-dimensional case with $n$ economic
observables, the monetary tangent space is $\Mtan_{s}
\subseteq \Real^{n}$, one coordinate per observable.
\end{remark}

\begin{proposition}[Monetary tangent space is a normed space]
\label{prop:normed-tangent}
Define the \emph{monetary norm} on $\Real^{n}$ by
\begin{equation}\label{eq:mnorm}
  \|v\|_{\Clock}
  := \|v\|_{2}
     \cdot \frac{1}{\sqrt{n}}\cdot\frac{1}{100},
  \quad v\in\Real^{n}.
\end{equation}
The pair $(\Real^{n}, \|\cdot\|_{\Clock})$ is a normed
vector space, and $\Mtan_{s}\subseteq\Real^{n}$ inherits
this norm.
\end{proposition}

\begin{proof}
The monetary norm is a scalar multiple of the Euclidean
norm, hence satisfies positivity, absolute homogeneity,
and the triangle inequality.  The prefactor
$(100\sqrt{n})^{-1}$ normalises the range so that the
maximum possible monetary tangent vector
$(100,\ldots,100)/100\sqrt{n}$ has monetary norm one.
\end{proof}

\begin{theorem}[Closure under composition]\label{thm:closure}
Let $f_{1},\ldots,f_{n}$ be S-transformed differentiable
quantities and let $\Phi:\Real^{n}\to\Real$ be a
differentiable function.  Then
$\tilde{h} := \Sval_{\mathcal{R}}(\Phi(f_{1},\ldots,f_{n}),
\Phi(f^{*}_{1},\ldots,f^{*}_{n}))$ satisfies
\begin{equation}
  \Mdiff{\tilde{h}}
  = \sum_{i=1}^{n}
    \frac{\partial\Phi}{\partial\tilde{f}_{i}}
    \cdot\Mdiff{\tilde{f}_{i}},
\end{equation}
and $\Mdiff{\tilde{h}}\in\Mtan_{s}$.
\end{theorem}

\begin{proof}
By the chain rule (Lemma~\ref{lem:chain}) applied to
$\tilde{h} = \Sval_{\mathcal{R}}(\Phi(\cdot),\cdot)$ and
the multivariate chain rule for $\Phi$:
$\dot{\tilde{h}} = \sum_{i}\partial\Phi/\partial\tilde{f}_{i}
\cdot\dot{\tilde{f}}_{i}$.  Dividing by $|\Gain|$ gives
the monetary derivative formula.  Since each
$\Mdiff{\tilde{f}_{i}}\in\Mtan_{s}$ and $\Mtan_{s}$ is
closed under finite linear combinations, the sum lies in
$\Mtan_{s}$.
\end{proof}

\begin{definition}[Monetary gradient]\label{def:mgradient}
For a scalar function $\Phi:\Real^{n}\to\Real$ of $n$
economic observables, the \emph{monetary gradient} is
\begin{equation}
  \nabla_{\Clock}\Phi
  := \left(\Mdiff{(\partial\Phi/\partial\tilde{f}_{i})}\right)_{i=1}^{n}
  \in \Mtan_{s}^{n}.
\end{equation}
\end{definition}

\begin{corollary}[Monetary gradient descent]\label{cor:mgd}
A quantity $\Phi$ evolves toward its truth value $\Phi^{*}$
in transaction time at rate
\begin{equation}
  \frac{d\tilde\Phi}{ds}(s)
  = -\|\nabla_{\Clock}\Phi\|_{\Clock}^{2}
    \cdot\frac{1}{\Mdiff{\tilde\Phi}}(\InvClock(s))
\end{equation}
along the monetary gradient flow
$\dot{\tilde{f}}_{i} = -\partial\Phi/\partial\tilde{f}_{i}$.
\end{corollary}

\begin{proof}
Along the gradient flow,
$d\tilde\Phi/dt = \sum_{i}(\partial\Phi/\partial\tilde{f}_{i})
\cdot\dot{\tilde{f}}_{i} = -\|\nabla\Phi\|^{2}$.  Dividing
by $|\Gain|$ gives $\Mdiff{\tilde\Phi} = -\|\nabla\Phi\|^{2}/|\Gain|$.
In transaction time this becomes
$d\tilde\Phi/ds = \Mdiff{\tilde\Phi}(\InvClock(s))$.
\end{proof}


% ============================================================
\section{Hierarchical Gear Networks and Telescoping
         Gear-Ratio Derivatives}
\label{sec:gear}
% ============================================================

We now develop the interaction between the monetary
derivative and the hierarchical gear network of nested
stopping times introduced in the context of portfolio
rebalancing.

\begin{definition}[Gear network]\label{def:gear}
A \emph{$K$-layer gear network} with thresholds
$\theta_{1}<\theta_{2}<\cdots<\theta_{K}$ is a sequence
of stopping times defined recursively as follows.  Set
$I_{0}(t) := \Gain(t,\tau)$ (the base-layer gain-loss).
For each $k=1,\ldots,K$:
\begin{enumerate}[nosep,label=(\roman*)]
\item Layer-$k$ transaction clock:
$\Clock_{k}(t) := \int_{0}^{t}|I_{k-1}(s)|\,ds$;
\item Layer-$k$ stopping time:
$\sigma_{k} := \inf\{t\geq 0 : \Clock_{k}(t) > \theta_{k}\}$;
\item Layer-$k$ accumulated imbalance:
$I_{k}(t) := \Clock_{k}(t) - \Clock_{k}(\sigma_{k}^{-})$
(reset to zero after each trigger).
\end{enumerate}
\end{definition}

\begin{remark}\label{rem:gear-hierarchy}
Definition~\ref{def:gear} defines a cascade: layer-1 fires
when cumulative gain-loss exceeds $\theta_{1}$; layer-2
fires when cumulative layer-1 activity exceeds $\theta_{2}$;
and so on.  The ratio $\theta_{k+1}/\theta_{k}$ is the
\emph{gear ratio} between layers $k$ and $k+1$.
\end{remark}

\begin{theorem}[Gear hierarchy mean crossing times]
\label{thm:gear-crossing}
Under the martingale difference structure and the
independence of increments across reset epochs,
the expected number of layer-$(k+1)$ firings in $[0,T]$
satisfies
\begin{equation}\label{eq:gear-frequency}
  \Expect[N_{k+1}(T)] = \frac{\theta_{k}}{\theta_{k+1}}
                         \cdot\Expect[N_{k}(T)]
                         \cdot(1+o(1))
\end{equation}
as $T\to\infty$, so that layer-$k$ fires approximately
$\theta_{k+1}/\theta_{k}$ times per layer-$(k+1)$ firing.
\end{theorem}

\begin{proof}
Between consecutive layer-$(k+1)$ firings, layer-$k$ must
accumulate activity $\theta_{k+1}$.  Each individual
layer-$k$ firing contributes exactly $\theta_{k}$ of
activity (by definition of the stopping time), so the
expected number of layer-$k$ firings required is
$\theta_{k+1}/\theta_{k}$.  By the renewal argument
\citep{Feller1971} applied to the i.i.d.\ inter-arrival
times of layer-$k$ firings (guaranteed by the reset
structure and the martingale difference property),
$\Expect[N_{k+1}(T)]\sim \Expect[N_{k}(T)]\cdot
(\theta_{k}/\theta_{k+1})$.
\end{proof}

\begin{definition}[Layer-$k$ monetary derivative]
\label{def:layered-mderiv}
For a quantity $f$ and layer $k\in\{1,\ldots,K\}$, the
\emph{layer-$k$ monetary derivative} is
\begin{equation}
  \MdiffK{f}{k}(t)
  := \frac{df/dt}{|\Clock_{k}(t)|/t}
  = \frac{t\,\dot{f}(t)}{|\Clock_{k}(t)|}.
\end{equation}
\end{definition}

\begin{theorem}[Telescoping gear-ratio formula]
\label{thm:telescope}
For any differentiable quantity $f$ and $K$-layer gear
network:
\begin{equation}\label{eq:telescope}
  \MdiffK{\tilde{f}}{K}(t)
  = \Mdiff{\tilde{f}}(t)
    \cdot \prod_{k=1}^{K-1} \rho_{k}(t),
\end{equation}
where $\rho_{k}(t) := \Clock_{k}(t)/\Clock_{k+1}(t)$
is the \emph{instantaneous gear ratio} at layer $k$.
\end{theorem}

\begin{proof}
Write $\MdiffK{\tilde{f}}{K} = \dot{\tilde{f}}/|\Clock_{K}/t|
= (\dot{\tilde{f}}/|\Clock_{1}/t|)\cdot
(|\Clock_{1}/t|/|\Clock_{K}/t|)$.
By telescoping,
$|\Clock_{1}/t|/|\Clock_{K}/t|
= \prod_{k=1}^{K-1}(|\Clock_{k}/t|/|\Clock_{k+1}/t|)
= \prod_{k=1}^{K-1}\rho_{k}(t)$.
Since $\Mdiff{\tilde{f}} = \dot{\tilde{f}}/|\Clock_{1}/t|$
(identifying the base layer with $k=1$), the result follows.
\end{proof}

\begin{corollary}[Gear amplification and attenuation]
\label{cor:gear-amp}
If thresholds satisfy $\theta_{1}<\theta_{2}<\cdots<\theta_{K}$,
then $\Clock_{k}(t)<\Clock_{k+1}(t)$ for all $t$ in the
long-run mean, so $\rho_{k}(t)<1$ for each $k$, and the
layer-$K$ monetary derivative is \emph{attenuated}
relative to the base layer:
$|\MdiffK{\tilde{f}}{K}| \leq |\Mdiff{\tilde{f}}|$.
Conversely, decreasing thresholds produce amplification.
\end{corollary}

\begin{proof}
By Theorem~\ref{thm:gear-crossing}, higher layers fire
less frequently, hence their transaction clocks advance
more slowly: $\Clock_{k+1}(t)\geq\Clock_{k}(t)$ in
expectation.  This gives $\rho_{k}(t)\leq 1$.
\end{proof}

\begin{theorem}[Gear-ratio stability under ergodic dynamics]
\label{thm:gear-stability}
Under a stationary ergodic gain-loss process with
$\Expect[|\Gain(t)|]=\bar{g}>0$, the instantaneous gear
ratios converge almost surely:
\begin{equation}
  \rho_{k}(t) \xrightarrow[t\to\infty]{a.s.}
  \bar{\rho}_{k} := \frac{\bar{g}_{k}}{\bar{g}_{k+1}},
\end{equation}
where $\bar{g}_{k} := \Expect[|I_{k-1}(t)|]$ is the
stationary mean activity of layer $k$.  Therefore,
the telescoping product converges to a deterministic limit
$\prod_{k=1}^{K-1}\bar{\rho}_{k}$.
\end{theorem}

\begin{proof}
By the ergodic theorem \citep{Birkhoff1931},
$\Clock_{k}(t)/t \to \bar{g}_{k}$ a.s.\ for each $k$.
Therefore $\rho_{k}(t) = \Clock_{k}(t)/\Clock_{k+1}(t)
\to \bar{g}_{k}/\bar{g}_{k+1} = \bar{\rho}_{k}$ a.s.
The product of $K-1$ almost surely convergent sequences
converges almost surely to their product.
\end{proof}


% ============================================================
\section{The No-Privileged-Level Theorem in Transaction Time}
\label{sec:npl}
% ============================================================

A central result of \citet{Sachikonye2024S} is the
\emph{Recursive No-Privileged-Level Theorem}: every
S-value decomposes into a triple of S-values at the next
finer depth, with the same algebraic laws holding at every
depth.  We now prove that this principle is preserved
exactly under the monetary derivative operation.

\begin{theorem}[No-Privileged-Level in transaction time]
\label{thm:npl}
Let $k,\ell\in\{1,\ldots,K\}$ be any two gear layers.
The algebraic laws of the monetary derivative at layer $k$
are identical to those at layer $\ell$: specifically,
\begin{enumerate}[nosep,label=(\roman*)]
\item \emph{Linearity:} $\MdiffK{(\alpha\tilde{f}+\beta\tilde{g})}{k}
      = \alpha\MdiffK{\tilde{f}}{k}+\beta\MdiffK{\tilde{g}}{k}$;
\item \emph{Product rule:}
      $\MdiffK{(\tilde{f}\tilde{g})}{k}
       = \tilde{g}\MdiffK{\tilde{f}}{k}
         +\tilde{f}\MdiffK{\tilde{g}}{k}$;
\item \emph{Chain rule:}
      $\MdiffK{(h\circ\tilde{f})}{k}
       = h'(\tilde{f})\MdiffK{\tilde{f}}{k}$;
\item \emph{Fundamental theorem:}
      $\tilde{f}(t_{2})-\tilde{f}(t_{1})
       = \int_{\Clock_{k}(t_{1})}^{\Clock_{k}(t_{2})}
         \MdiffK{\tilde{f}}{k}(\InvClock_{k}(s))\,ds$.
\end{enumerate}
These laws hold identically for every $k\in\{1,\ldots,K\}$.
\end{theorem}

\begin{proof}
Each law (i)--(iv) is proved by dividing the corresponding
calendar-time identity by $|\Clock_{k}(t)|/t$, which
appears as a common factor and cancels.  The resulting
identities are identical in form for all $k$, since the
specific layer enters only through the denominator
$\Clock_{k}$, and all the algebraic manipulations
(linearity, Leibniz rule, chain rule, FTC) proceed
identically regardless of which layer's clock is used.
\end{proof}

\begin{corollary}[Level-independent calculus]\label{cor:level-indep}
The set of monetary derivatives $\{\MdiffK{\tilde{f}}{k}\}_{k=1}^{K}$
forms a family of calculi indexed by gear layer, all
sharing the same algebraic structure.  No layer carries
greater algebraic privilege than any other.
\end{corollary}

\begin{remark}\label{rem:npl-interpretation}
Theorem~\ref{thm:npl} is the analogue in transaction time
of the Recursive No-Privileged-Level Theorem of
\citet{Sachikonye2024S}: just as every S-value is its
own sub-problem at the next depth (with the same laws),
every gear layer defines an equivalent calculus of monetary
derivatives.  The gear ratios $\rho_{k}$ encode the
\emph{quantitative} differences between layers (the
amplification or attenuation of rate signals) without
introducing any \emph{qualitative} algebraic asymmetry.
\end{remark}

\begin{theorem}[S-entropy floor preservation across layers]
\label{thm:floor-layers}
For every bounded receiver $\mathcal{R}$ with floor
$\Sfloor(\mathcal{R})>0$ and every gear layer $k$:
\begin{equation}
  \tilde{f}(t) \;\in\;
  [\Sfloor(\mathcal{R}), 100] \quad \forall\, t\geq 0,
\end{equation}
and this constraint propagates layer by layer:
$\MdiffK{\tilde{f}}{k}$ is bounded by
$(100-\Sfloor(\mathcal{R}))/\Clock_{k,\min}$
whenever $\Clock_{k}(t)\geq\Clock_{k,\min}>0$.
\end{theorem}

\begin{proof}
$\tilde{f}(t)\in[\Sfloor(\mathcal{R}),100]$ by
Axioms (S1)--(S3) of the S-functional.  Therefore
$|\dot{\tilde{f}}(t)|\leq 100-\Sfloor(\mathcal{R})$
almost everywhere, since $\tilde{f}$ cannot leave the
bounded interval.  Dividing by $|\Clock_{k}(t)|/t
\geq \Clock_{k,\min}/t$ gives
$|\MdiffK{\tilde{f}}{k}|\leq
t(100-\Sfloor)/|\Clock_{k}(t)|
\leq (100-\Sfloor)/\Clock_{k,\min}$.
\end{proof}


% ============================================================
\section{Ergodic Theory of Monetary Derivatives}
\label{sec:ergodic}
% ============================================================

\begin{theorem}[Ergodic consistency of monetary derivatives]
\label{thm:ergodic}
Let the return process $(R(t,\tau))_{t\geq 0}$ be
stationary and ergodic under $\Prob$, with invariant
measure $\pi$.  Let $f$ be a measurable function of
the returns with $\Expect_{\pi}[|\dot{\tilde{f}}|/\bar{g}]<\infty$.
Then
\begin{equation}\label{eq:ergodic-mderiv}
  \frac{1}{S}\int_{0}^{S}
  \Mdiff{\tilde{f}}(\InvClock(s))\,ds
  \xrightarrow[S\to\infty]{a.s.}
  \Expect_{\pi}\!\left[\Mdiff{\tilde{f}}\right]
  = \frac{\Expect_{\pi}[\dot{\tilde{f}}]}{\bar{g}}.
\end{equation}
\end{theorem}

\begin{proof}
By Proposition~\ref{prop:mderiv-transtime},
$\Mdiff{\tilde{f}}(\InvClock(s)) = dY/ds(s)$ where
$Y(s)=\tilde{f}(\InvClock(s))$.  The time-average on the
left of \eqref{eq:ergodic-mderiv} is the Birkhoff average
of $dY/ds$ in transaction time.  Since the return process
is ergodic and $\Clock$ subordinates the process to the
transaction time, $Y(s)$ inherits ergodicity from the
underlying stationary process (see \citealt{Revuz1999},
Theorem~4.1.2).  Birkhoff's ergodic theorem then gives
$S^{-1}\int_{0}^{S}(dY/ds)\,ds\to\Expect_{\pi}[dY/ds]$ a.s.
The expectation equals $\Expect_{\pi}[\dot{\tilde{f}}]/\bar{g}$
by the ergodic theorem applied to $\dot{\tilde{f}}$ in
calendar time, combined with $\Clock(t)/t\to\bar{g}$ a.s.
\end{proof}

\begin{theorem}[Convergence rate of the Birkhoff average]
\label{thm:clt-mderiv}
Under the conditions of Theorem~\ref{thm:ergodic} and
additionally assuming $\Expect_{\pi}[(\dot{\tilde{f}}/\bar{g})^{2}]
<\infty$ and a finite mixing time:
\begin{equation}\label{eq:clt-rate}
  \left\|
    \frac{1}{S}\int_{0}^{S}\Mdiff{\tilde{f}}(\InvClock(s))\,ds
    - \Expect_{\pi}\!\left[\Mdiff{\tilde{f}}\right]
  \right\|_{2}
  = O(S^{-1/2}).
\end{equation}
\end{theorem}

\begin{proof}
Under the mixing condition, the central limit theorem for
ergodic processes \citep{Billingsley1999} applies to the
partial sums $\sum_{k=1}^{K}(dY/ds)|_{s=k\Delta}$.  In
the continuous-time limit this gives the $O(S^{-1/2})$
rate of convergence, consistent with the Birkhoff theorem
applied to a square-integrable function.
\end{proof}

\begin{corollary}[Monetary law of large numbers]\label{cor:mlln}
For a discrete sequence of transaction-time observations
$s_{1}<s_{2}<\cdots<s_{K}$ with spacing $\Delta s$:
\begin{equation}
  \frac{1}{K}\sum_{k=1}^{K}
  \Mdiff{\tilde{f}}(\InvClock(s_{k}))
  \xrightarrow[K\to\infty]{a.s.}
  \Expect_{\pi}\!\left[\Mdiff{\tilde{f}}\right].
\end{equation}
\end{corollary}

\begin{proof}
Apply Theorem~\ref{thm:ergodic} to the discrete-time
Birkhoff average; ergodicity of the discrete skeleton
follows from that of the continuous process
\citep{MeynTweedie1993}.
\end{proof}

\begin{theorem}[Inter-layer ergodic consistency]
\label{thm:inter-layer-ergodic}
Under the same ergodicity conditions, for any two layers
$j,k\in\{1,\ldots,K\}$:
\begin{equation}
  \frac{\Expect_{\pi}[\MdiffK{\tilde{f}}{k}]}
       {\Expect_{\pi}[\MdiffK{\tilde{f}}{j}]}
  = \frac{\bar{g}_{j}}{\bar{g}_{k}},
\end{equation}
and the ratio is determined solely by the mean activity
rates $\bar{g}_{j},\bar{g}_{k}$ of the respective layers.
\end{theorem}

\begin{proof}
$\Expect_{\pi}[\MdiffK{\tilde{f}}{k}]
= \Expect_{\pi}[\dot{\tilde{f}}]/\bar{g}_{k}$ by
Theorem~\ref{thm:ergodic} applied at layer $k$, and
similarly for layer $j$.  Dividing gives the ratio.
\end{proof}


% ============================================================
\section{Worked Examples}
\label{sec:examples}
% ============================================================

\subsection{Price velocity in transaction time}

Let $P(t)$ be the mid-price of a single asset with truth
value $P^{*}(t)=$ fundamental value.  The S-transform is
$\tilde{P}(t) = \Sval_{\mathcal{R}}(P(t),P^{*}(t))$.
Suppose prices move at rate $\dot{P}(t)$ and the portfolio
generates gain-loss $\Gain(t)$.  Then the monetary price
velocity is
\begin{equation}
  \Mdiff{\tilde{P}}(t)
  = \frac{d\Sval(P,P^{*})/dt}{|\Gain(t)|}
  = \frac{\partial\Sval/\partial P \cdot \dot{P}}{|\Gain|}.
\end{equation}
This answers the question: how fast is the price diverging
from fundamental per unit of economic activity?  A large
$\Mdiff{\tilde{P}}$ means rapid mispricing per transaction
unit; a small value means the mispricing is moving slowly
relative to the economic intensity of the market.

\subsection{Cross-asset comparison: price and volume}

Consider two assets with price $P$ and volume $V$.  Their
individual monetary derivatives $\Mdiff{\tilde{P}}$ and
$\Mdiff{\tilde{V}}$ are both in S-units per transaction
unit.  Their sum
\begin{equation}
  \Mdiff{\tilde{P}} + \Mdiff{\tilde{V}}
\end{equation}
is a single dimensionless number representing the
\emph{total rate of informational divergence} of both
price and volume from their respective truths, per unit
of transactional activity.  This would be meaningless in
calendar time (dollars-per-second plus shares-per-second
is ill-defined); in transaction time with S-entropy
normalisation, it is well-defined by Corollary~\ref{cor:additivity}.

\subsection{Two-layer gear derivative of volatility}

Let $\sigma(t)^{2}$ be the realised variance of a
portfolio.  Using a two-layer gear network with thresholds
$\theta_{1}=0.1$ and $\theta_{2}=0.5$:
\begin{equation}
  \MdiffK{\tilde{\sigma}}{2}(t)
  = \Mdiff{\tilde{\sigma}}(t) \cdot \rho_{1}(t)
  = \frac{\dot{\tilde\sigma}}{|\Gain|}
    \cdot \frac{\Clock_{1}(t)}{\Clock_{2}(t)}.
\end{equation}
The factor $\rho_{1}(t) = \Clock_{1}(t)/\Clock_{2}(t)<1$
attenuates the base-layer signal: the layer-2 monetary
derivative of volatility represents the rate at which
volatility deviates from truth \emph{only as reckoned by
the slow macro-rebalancing clock}, filtering out
high-frequency noise from the base-layer activity.


% ============================================================
\section{Implications and Open Questions}
\label{sec:conclude}
% ============================================================

\subsection{Relation to dimensional analysis}

The monetary derivative in S-units is the first
\emph{automatically dimensionless} rate operator for
economic systems.  Classical dimensional analysis
\citep{Bridgman1922} requires explicit unit-tracking;
the Buckingham $\Pi$ theorem \citep{Buckingham1914}
constructs dimensionless groups post-hoc.  The monetary
derivative achieves dimensionlessness directly, by using
S-entropy as the universal scalar field and accumulated
absolute gain-loss as the universal clock.

\subsection{Connection to information geometry}

The monetary tangent space $\Mtan_{s}$ with the monetary
norm $\|\cdot\|_{\Clock}$ is a Riemannian manifold when
the S-functional varies smoothly with the underlying
process.  The monetary gradient of
Definition~\ref{def:mgradient} defines a natural gradient
flow on this manifold.  The connection to information
geometry \citep{AmariFTJ1987} and the Fisher information
metric is an open question.

\subsection{Stochastic monetary calculus}

When $\Gain(t)$ is stochastic, the monetary derivative
$\Mdiff{f} = \dot{f}/|\Gain|$ is a random variable.  A
full stochastic calculus in transaction time---including
It\^o's formula for $d\tilde{f}$ in $d\Clock$---follows
from the theory of stochastic time changes
\citep{KaratzasShreve1991}, but the S-entropy nonlinearity
requires careful treatment.  We leave the development of
a complete It\^o calculus for S-entropy-valued processes
as an open question.

\subsection{Gear networks with variable thresholds}

Theorem~\ref{thm:telescope} and the No-Privileged-Level
result (Theorem~\ref{thm:npl}) assume fixed thresholds
$\theta_{k}$.  Adaptive threshold sequences
$\theta_{k}(t)$ that respond to realised volatility would
make the gear ratios $\rho_{k}(t)$ adaptive as well.
Whether the No-Privileged-Level structure is preserved
under such adaptations is an open question with practical
significance for adaptive rebalancing algorithms.

\subsection{Summary}

We have constructed a self-consistent calculus in which:
\begin{enumerate}[nosep]
\item the fundamental rate operator is the monetary
      derivative $d/d\Clock$, measuring change per unit
      of accumulated gain-loss intensity;
\item S-entropy provides a universal, dimensionless scalar
      field that makes rates of change of heterogeneous
      quantities addable;
\item the monetary tangent space is a normed vector space
      closed under the monetary derivative;
\item a telescoping gear-ratio formula propagates
      monetary derivatives through any finite gear network;
\item the No-Privileged-Level theorem ensures that the
      same algebraic laws hold at every gear layer; and
\item under stationary ergodic dynamics, time-averaged
      monetary derivatives converge almost surely at rate
      $O(K^{-1/2})$.
\end{enumerate}
The calculus reduces to standard calculus when $|\Gain|$
is constant, and subsumes Clark's subordination framework
as the special case where S-entropy is replaced by the
identity function.

\vspace{6pt}
\bibliographystyle{plainnat}
\bibliography{references}

\end{document}
